\documentclass[10pt,a4paper]{amsart}
\usepackage[margin=19mm]{geometry}
\usepackage{amsmath,amssymb,mathtools}
\usepackage[scr=rsfs,cal=euler]{mathalfa}
\usepackage{xfrac}
\usepackage[unicode,hidelinks,bookmarks=false]{hyperref}
\newcommand{\ie}{i.e.\ }
\newcommand{\cf}{cf.\ }
\newcommand{\resp}{respectively\ }
\newcommand{\Subsection}{\subsection}
\providecommand{\nobreakdash}{\nobreak\hbox{-}}
\setlength{\emergencystretch}{2em}
\multlinegap0pt
\usepackage{amssymb}
\usepackage{caption}
\usepackage{booktabs}
\usepackage{float}
\newtheorem{theorem}{Theorem}
\newcommand{\ptop}{p_{\top}}
\title{Godel Implication on Finite Chains: Truth Tables and Catalan-Bracketing Enumerations}
\author{Volkan Yildiz}
\date{Source: arXiv:2602.16135v1 (CC BY 4.0)}
\begin{document}
\maketitle

\noindent\emph{Source and licence.} Godel Implication on Finite Chains: Truth Tables and Catalan-Bracketing Enumerations. Version arXiv:2602.16135v1 (CC BY 4.0), distributed by arXiv under the Creative Commons Attribution 4.0 International licence, http://creativecommons.org/licenses/by/4.0/, as recorded in the arXiv licence metadata for this exact version and shown on the abstract page https://arxiv.org/abs/2602.16135v1. Full licence notice: \texttt{licenses/arxiv-2602.16135-CC-BY-4.0.txt}.\par\smallskip

\noindent\textit{Editorial scope.} Counted unit: the article's theorem thm:p1limit (source0.tex lines 1038--1049). The article's complete original proof of it is delivered with it (source0.tex lines 1051--1176, one \texttt{proof} environment of 126 lines). No mathematics is written, repaired or re-derived: the statement and the proof are the article's own lines, in the article's own notation, and no retained line is substituted or re-flowed.  No further source-local context is needed: every cross-reference of every retained line resolves inside the two delivered spans.  Nothing was omitted from any retained span, so the package carries no omission record.  No retained line contains a citation, so the wrapper carries no bibliography and no bibliography entry.  No retained line contains a figure environment, an image-inclusion command or drawing code, so no image is delivered and the package declares no asset dependency.  Editorial formatting only, in addition: an AMS \texttt{amsart} preamble that restates the article's own declarations and macros used by the retained lines, namely the environments \texttt{theorem} on separate counters, together with the standard packages those lines need.  The retained environments print in this document as Theorem 1 in order of appearance, whereas the article prints the same items with the numbering of its own full document.  The complete native source of the article as distributed in the arXiv e-print for this exact version is bundled unchanged as \texttt{sources/194850/source0.tex}.
\begin{theorem}[As $m\to\infty$]
\label{thm:p1limit}
Let \(r=\frac1{4m}\) and let \((w_k)_{k\ge0}\) be defined by \(w_{k+1}=\Psi_r(w_k)\) with
\(w_0=\frac12\). Define
\[
\ptop(m):=\prod_{k=0}^{m-2}\partial_v\Psi_r(w_k).
\]
Then
\[
\boxed{\;\;\ptop(m)\longrightarrow \frac{1}{\sqrt{3}}.\;\;}
\]
\end{theorem}

\begin{proof}
Write \(u_k:=1-w_k\). In our situation \(\Psi_r\) admits the closed form
\[
\Psi_r(w)=\frac{1+w-\sqrt{(1-w)^2+\frac1m}}{2},
\qquad\text{so that}\qquad
u_{k+1}=\frac{u_k+\sqrt{u_k^2+\frac1m}}{2},
\]
and
\[
\partial_v\Psi_r(w)
=\frac12\left(1+\frac{u}{\sqrt{u^2+\frac1m}}\right),
\qquad u=1-w.
\]
Since \(u_{k+1}-u_k=\bigl(\sqrt{u_k^2+\frac1m}-u_k\bigr)/2\ge0\), the sequence \((u_k)\) is increasing.
Moreover \(u_0=\frac12\), and if \(u_k\le 1\) then
\[
u_{k+1}\le \frac{u_k+\sqrt{u_k^2+1}}{2}\le \frac{u_k+(u_k+1)}{2}\le 1.
\]
Hence
\[
\frac12\le u_k\le 1\qquad (0\le k\le m).
\]

\medskip
\noindent
\emph{Uniform expansion of \(\log(\partial_v\Psi_r(w_k))\).}
Let \(\varepsilon:=\frac{1}{m u^2}\). For \(u\in[1/2,1]\) we have \(u^2\ge 1/4\), hence
\[
0\le \varepsilon=\frac{1}{m u^2}\le \frac{4}{m},
\]
so \(\varepsilon\to0\) uniformly as \(m\to\infty\). Using
\[
\frac{u}{\sqrt{u^2+\frac1m}}=\frac{1}{\sqrt{1+\varepsilon}}
=1-\frac{\varepsilon}{2}+O(\varepsilon^2),
\]
we obtain, uniformly for \(u\in[1/2,1]\),
\begin{equation}
\label{eq:derivexp_rewrite}
\partial_v\Psi_r(1-u)
=1-\frac{1}{4m u^2}+O\!\left(\frac{1}{m^2}\right).
\end{equation}
Set
\[
x_k:=\frac{1}{4m u_k^2},\qquad \delta_k:=\partial_v\Psi_r(w_k)-\bigl(1-x_k\bigr).
\]
Then \(0\le x_k\le 1/m\) and, by \eqref{eq:derivexp_rewrite}, \(|\delta_k|\le C/m^2\) uniformly for \(k\le m\).
In particular \(x_k-\delta_k=O(1/m)\) uniformly, hence for \(m\) large enough \(\partial_v\Psi_r(w_k)\in(1/2,1)\)
uniformly and we may use \(\log(1-y)=-y+O(y^2)\) uniformly:
\[
\log\bigl(\partial_v\Psi_r(w_k)\bigr)
=\log\bigl(1-(x_k-\delta_k)\bigr)
=-(x_k-\delta_k)+O\bigl((x_k-\delta_k)^2\bigr).
\]
Summing from \(k=0\) to \(m-2\) and using
\[
\sum_{k\le m}x_k^2\le m\cdot (1/m^2)=O(1/m),\qquad
\sum_{k\le m}|\delta_k|=O(1/m),
\]
we obtain
\begin{equation}
\label{eq:logp1_rewrite}
\log \ptop(m)
=\sum_{k=0}^{m-2}\log\bigl(\partial_v\Psi_r(w_k)\bigr)
=-\sum_{k=0}^{m-2}\frac{1}{4m u_k^2}+o(1).
\end{equation}

\medskip
\noindent
\emph{Euler approximation.}
Let \(h:=1/m\) and \(t_k:=kh\). Let \(u(t)\) be the solution of
\[
u'(t)=\frac{1}{4u(t)},\qquad u(0)=\frac12,
\]
so \(u(t)^2=\frac14+\frac{t}{2}\) and \(u(t)\in[1/2,1]\) for \(t\in[0,1]\).
From the exact recurrence,
\[
u_{k+1}-u_k=\frac{\sqrt{u_k^2+h}-u_k}{2}
=\frac{h}{2\bigl(\sqrt{u_k^2+h}+u_k\bigr)}
=\frac{h}{4u_k}+O(h^2),
\]
uniformly for \(u_k\in[1/2,1]\). Thus
\[
u_{k+1}=u_k+h f(u_k)+\rho_k,\qquad f(u)=\frac{1}{4u},\qquad |\rho_k|\le C_1 h^2.
\]
By Taylor's theorem applied to the ODE solution,
\[
u(t_{k+1})=u(t_k)+h f\bigl(u(t_k)\bigr)+\eta_k,\qquad |\eta_k|\le C_2 h^2,
\]
uniformly for \(t_k\in[0,1]\). Setting \(e_k:=u_k-u(t_k)\) and subtracting gives
\[
e_{k+1}=e_k+h\bigl(f(u_k)-f(u(t_k))\bigr)+(\rho_k-\eta_k).
\]
Since \(|f'(u)|=1/(4u^2)\le 1\) on \([1/2,1]\), \(f\) is Lipschitz with constant \(L=1\), hence
\[
|e_{k+1}|\le (1+Lh)|e_k|+C_3 h^2.
\]
Iterating and using \(e_0=0\) yields \(|e_k|\le C_4 h\) uniformly for \(0\le k\le m\), i.e.
\[
\max_{0\le k\le m}|u_k-u(t_k)|=O(h)=O(1/m).
\]

\medskip
\noindent
\emph{Riemann-sum limit.}
Since \(t\mapsto 1/u(t)^2\) is continuous on \([0,1]\) and bounded, and \(u_k=u(t_k)+O(1/m)\) uniformly, we have
\[
\frac{1}{m}\sum_{k=0}^{m-2}\frac{1}{u_k^2}
=\frac{1}{m}\sum_{k=0}^{m-2}\frac{1}{u(t_k)^2}+o(1)
\longrightarrow \int_0^1\frac{dt}{u(t)^2}.
\]
Combining this with \eqref{eq:logp1_rewrite} gives
\[
\log \ptop(m)\longrightarrow -\int_0^1\frac{dt}{4u(t)^2}.
\]
Using \(u(t)^2=\frac14+\frac{t}{2}\),
\[
\int_0^1\frac{dt}{4u(t)^2}
=\int_0^1\frac{dt}{1+2t}
=\frac12\log 3.
\]
Hence \(\log \ptop(m)\to -\frac12\log 3\), i.e.
\[
\ptop(m)\to e^{-(1/2)\log 3}=\frac{1}{\sqrt3},
\]
as claimed.
\end{proof}

\end{document}
